\section{Irreducible components of spectra}
\label{section-irreducible}

\noindent
We show that irreducible components of
the spectrum of a ring correspond to the
minimal primes in the ring.

\begin{lemma}
\label{lemma-irreducible}
Let $R$ be a ring.
\begin{enumerate}
\item For a prime $\mathfrak p \subset R$ the closure
of $\{\mathfrak p\}$ in the Zariski topology is $V(\mathfrak p)$.
In a formula $\overline{\{\mathfrak p\}} = V(\mathfrak p)$.
\item The irreducible closed subsets of $\Spec(R)$ are
exactly the subsets $V(\mathfrak p)$, with $\mathfrak p \subset R$
a prime.
\item The irreducible components (see Topology,
Definition \ref{topology-definition-irreducible-components})
of $\Spec(R)$ are  exactly the subsets $V(\mathfrak p)$,
with $\mathfrak p \subset R$ a minimal prime.
\end{enumerate}
\end{lemma}

\begin{proof}
Note that if $ \mathfrak p \in V(I)$, then
$I \subset \mathfrak p$. Hence,
clearly $\overline{\{\mathfrak p\}} = V(\mathfrak p)$.
In particular $V(\mathfrak p)$ is the closure of
a singleton and hence irreducible.
The second assertion implies the third.
To show the second, let
$V(I) \subset \Spec(R)$ with $I$ a radical ideal.
If $I$ is not prime, then choose $a, b\in R$, $a, b\not \in I$
with $ab\in I$. In this case $V(I, a) \cup V(I, b) = V(I)$,
but neither $V(I, b) = V(I)$ nor $V(I, a) = V(I)$, by
Lemma \ref{lemma-Zariski-topology}. Hence $V(I)$ is not
irreducible.
\end{proof}

\noindent
In other words, this lemma shows that every irreducible closed
subset of $\Spec(R)$ is of the form $V(\mathfrak p)$ for
some prime $\mathfrak p$. Since $V(\mathfrak p) = \overline{\{\mathfrak p\}}$
we see that each irreducible closed subset has a unique generic point,
see Topology, Definition \ref{topology-definition-generic-point}.
In particular, $\Spec(R)$ is a sober topological space.
We record this fact in the following lemma.

\begin{lemma}
\label{lemma-spec-spectral}
The spectrum of a ring is a spectral space, see Topology, Definition
\ref{topology-definition-spectral-space}.
\end{lemma}

\begin{proof}
Formally this follows from Lemma \ref{lemma-irreducible} and
Lemma \ref{lemma-topology-spec}. See also discussion above.
\end{proof}

\begin{lemma}
\label{lemma-irreducible-components-containing-x}
Let $R$ be a ring. Let $\mathfrak p \subset R$ be a prime.
\begin{enumerate}
\item the set of irreducible closed subsets of $\Spec(R)$
passing through $\mathfrak p$ is in one-to-one correspondence with
primes $\mathfrak q \subset R_{\mathfrak p}$.
\item The set of irreducible components of $\Spec(R)$ passing through
$\mathfrak p$ is in one-to-one correspondence with minimal
primes $\mathfrak q \subset R_{\mathfrak p}$.
\end{enumerate}
\end{lemma}

\begin{proof}
Follows from Lemma \ref{lemma-irreducible}
and the description of $\Spec(R_\mathfrak p)$ in
Lemma \ref{lemma-spec-localization} which shows that
$\Spec(R_\mathfrak p)$ corresponds to primes $\mathfrak q$ in $R$
with $\mathfrak q \subset \mathfrak p$.
\end{proof}

\begin{lemma}
\label{lemma-standard-open-containing-maximal-point}
Let $R$ be a ring.
Let $\mathfrak p$ be a minimal prime of $R$.
Let $W \subset \Spec(R)$ be a quasi-compact open
not containing the point $\mathfrak p$. Then there
exists an $f \in R$, $f \not \in \mathfrak p$ such
that $D(f) \cap W = \emptyset$.
\end{lemma}

\begin{proof}
Since $W$ is quasi-compact we may write it as a finite union
of standard affine opens $D(g_i)$, $i = 1, \ldots, n$.
Since $\mathfrak p \not \in W$ we have $g_i \in \mathfrak p$ for
all $i$. By Lemma \ref{lemma-minimal-prime-reduced-ring}
each $g_i$ is nilpotent in $R_{\mathfrak p}$. Hence we can find
an $f \in R$, $f \not \in \mathfrak p$ such that for all $i$ we have
$f g_i^{n_i} = 0$ for some $n_i > 0$. Then $D(f)$ works.
\end{proof}

\begin{lemma}
\label{lemma-ring-with-only-minimal-primes}
Let $R$ be a ring. Let $X = \Spec(R)$ as a topological space.
The following are equivalent
\begin{enumerate}
\item $X$ is profinite,
\item $X$ is Hausdorff,
\item $X$ is totally disconnected.
\item every quasi-compact open of $X$ is closed,
\item there are no nontrivial inclusions between its prime ideals,
\item every prime ideal is a maximal ideal,
\item every prime ideal is minimal,
\item every standard open $D(f) \subset X$ is closed, and
\item add more here.
\end{enumerate}
\end{lemma}

\begin{proof}
First proof. It is clear that (5), (6), and (7) are equivalent.
It is clear that (4) and (8) are equivalent as every quasi-compact
open is a finite union of standard opens.
The implication (7) $\Rightarrow$ (4) follows from
Lemma \ref{lemma-standard-open-containing-maximal-point}.
Assume (4) holds. Let $\mathfrak p, \mathfrak p'$ be distinct
primes of $R$. Choose an $f \in \mathfrak p'$, $f \not \in \mathfrak p$
(if needed switch $\mathfrak p$ with $\mathfrak p'$).
Then $\mathfrak p' \not \in D(f)$ and $\mathfrak p \in D(f)$.
By (4) the open $D(f)$ is also closed.
Hence $\mathfrak p$ and $\mathfrak p'$ are in disjoint open
neighbourhoods whose union is $X$. Thus $X$ is Hausdorff and totally
disconnected. Thus (4) $\Rightarrow$ (2) and (3).
If (3) holds then there cannot be any specializations
between points of $\Spec(R)$ and we see that (5) holds.
If $X$ is Hausdorff then every point is closed, so (2) implies (6).
Thus (2), (3), (4), (5), (6), (7) and (8) are equivalent.
Any profinite space is Hausdorff, so (1) implies (2).
If $X$ satisfies (2) and (3), then $X$ (being quasi-compact by
Lemma \ref{lemma-quasi-compact}) is profinite by
Topology, Lemma \ref{topology-lemma-profinite}.

\medskip\noindent
Second proof. Besides the equivalence of (4) and (8) this follows
from Lemma \ref{lemma-spec-spectral} and purely topological facts, see
Topology, Lemma \ref{topology-lemma-characterize-profinite-spectral}.
\end{proof}













